\SourceLeaf{287}%
all the closer to the value $360^0$ the smaller $\omega, \omega'$ are. I.\,e., the tangent to the curve has changed its inclination by $360^0$ from beginning to end.

If the curve intersects itself once, these conclusions are modified. If the endpoints of $s = \sigma$, $s' = \sigma'$ coincide, then $(s', s)$ undergoes a jump at these values. Namely, with $s = \sigma$ held constant, $(s', s)$ changes all at once as $s'$ reaches the value $\sigma'$, by $\pm 180^0$. To decide which value applies here, one must consider the values for $s = \sigma \mp \omega$, where $\omega$ is infinitesimal. One sees that as $s' - \sigma'$ goes from a [negative] value [onward] through ever smaller negative values and grows again through $0$, $(s', \sigma - \omega)$ increases or decreases by an amount that differs infinitesimally from $180^0$, according as at $s'$ one passes from the inner to the outer side of the line $s'$ or vice versa. For $(s', \sigma + \omega)$ it is the other way round. Thus in general for a fixed $s'$ greater than $\sigma'$, $(s', s)$ on passing through the value $s = \sigma$ suddenly decreases (increases) by $360^0$, and the change in $(s', \sigma)$, as $s'$ passes through the value $\sigma'$, by $\pm 180^0$ actually holds with both signs at once, according as one regards $s = \sigma$ as the last of the values taken \textls{up to} $\sigma$ or as the first of the values to be taken onward \textls{from} $\sigma$. The solid here has the shape of a spiral staircase. The change in the direction of the tangent now becomes either $720^0$ or $0^0$, but one can dispense with this consideration altogether by restricting the present method to the case in which the curve does \textls{not} intersect itself.

\begin{center}
[8.]
\end{center}

In general, much in this theory can be greatly simplified if one assumes that polygons can be decomposed into triangles. This can be proved as follows.

First it is clear that every figure can be divided by a single dividing line (straight, shortest) into two figures that together have one fewer inward angle. Repeating this kind of division thus splits the polygon into others that have no inward angles at all.

Now let polygon $P$ have only outward angles at $(1), (2), (3) .\,.\,. (n)$. I claim that the line $(1)(3)$ will lie entirely inside the polygon.
